Modern classifiers are often selected by in-distribution accuracy, yet
downstream reliability can depend on geometric properties of the learned
representation that accuracy does not reveal. This is especially important for
feature-based uncertainty and out-of-distribution (OOD) detectors, which read
distances, norms, covariances, and local neighborhoods in the penultimate
feature space. We study whether optimizer choice preserves the geometry required
by these detector families. Rather than proposing a new detector, we use SGD,
Adam, and AdamW as controlled interventions on representation geometry and
evaluate how their optimizer-induced feature profiles are read by logit-based
and feature-based OOD scores. In accuracy-controlled selected configurations on
CIFAR-10 with WideResNet-28-10, we find that models with comparable ID
performance can exhibit distinct penultimate geometry fingerprints, including
differences in neural-collapse metrics, feature norms, within-class dispersion,
and class separation. These differences are not uniformly reflected in detector
behavior: logit-based scores can remain competitive while raw Mahalanobis and
kNN scores degrade substantially in AdamW-selected regimes. Post-hoc L2
normalization recovers much of this raw feature-detector gap, indicating that
radial or norm-coupled geometry is a major failure channel, while
covariance-side controls show more limited and detector-dependent effects. Our
results therefore do not support an optimizer-good/bad ranking. They support a
geometry--detector compatibility view: optimization improvements may preserve,
distort, or mask the representation geometry required by a downstream
uncertainty estimator. This controlled study suggests that optimizer choice
should be evaluated jointly with penultimate geometry and detector-family
readout, not by ID accuracy alone.
